(1) 设转动矢量 \(\boldsymbol\Omega=\Omega\hat{\boldsymbol z}\)，惯性系位置与旋转系分量的几何关系为 \(\boldsymbol r=R(\Omega t)\boldsymbol r'\)。两次对时间求导得
\[
\boldsymbol a_{\rm in}=R(\Omega t)[\ddot{\boldsymbol r}'+2\boldsymbol\Omega\times\dot{\boldsymbol r}'+\boldsymbol\Omega\times(\boldsymbol\Omega\times\boldsymbol r')].
\]
由于弹簧力是中心力，旋转系满足
\[
\ddot{\boldsymbol r}'+2\boldsymbol\Omega\times\dot{\boldsymbol r}'+(\omega_0^2-\Omega^2)\boldsymbol r'=0.
\]
因此科里奥利力为 \(-2m\boldsymbol\Omega\times\dot{\boldsymbol r}'\)、离心力为 \(+m\Omega^2\boldsymbol r'\)。

(2) 令 \(z'=x'+iy'\)，平面叉乘 \(\hat z\times\boldsymbol r'\) 对应乘以 \(i\)，故
\[
\ddot z'+2i\Omega\dot z'+(\omega_0^2-\Omega^2)z'=0.
\]
取解 \(e^{i\nu t}\)，特征方程 \(-\nu^2-2\Omega\nu+\omega_0^2-\Omega^2=0\)，于是 \(\nu=-\Omega\pm\omega_0\)。因此
\[
z'(t)=A e^{i(\omega_0-\Omega)t}+B e^{-i(\omega_0+\Omega)t}.
\]
常数 \(A,B\) 为任意复数，代表相反圆极化及不同复振幅；两个有符号频率不能仅取绝对值后丢掉传播方向。

(3) 在 \(t=0\) 旋转系初位置 \(z'(0)=a\)；旋转坐标速度与惯性速度之差为 \(-\boldsymbol\Omega\times\boldsymbol r\)，故 \(\dot z'(0)=i(v-\Omega a)\)。联立得
\[
A=\frac12\left(a+\frac{v}{\omega_0}\right),\qquad
B=\frac12\left(a-\frac{v}{\omega_0}\right).
\]
乘回 \(e^{i\Omega t}\) 恢复惯性系
\[
z(t)=a\cos(\omega_0t)+i\frac v{\omega_0}\sin(\omega_0t).
\]
所以真实轨道是以原点为中心的椭圆（退化参数允许直线或圆），不因观察者旋转改变。

(4) 取旋转系速度平方 \(|\dot{\boldsymbol r}'|^2\)，对运动方程点乘 \(\dot{\boldsymbol r}'\)，科里奥利项始终不做功，故
\[
E_J=\frac m2|\dot{\boldsymbol r}'|^2+\frac m2(\omega_0^2-\Omega^2)|\boldsymbol r'|^2
\]
守恒。利用 \(\boldsymbol v_{\rm in}=R(\dot{\boldsymbol r}'+\boldsymbol\Omega\times\boldsymbol r')\) 展开，可得
\[
E_J=E_{\rm in}-\Omega L_z,\qquad
E_{\rm in}=\tfrac m2|\boldsymbol v_{\rm in}|^2+\tfrac k2r^2.
\]
二者因中心力分别守恒。在 \(\Omega=\omega_0\) 时一个旋转模态频率为零，仅说明该圆运动在共转坐标中静止；惯性系弹簧势仍正定且全部运动有界，不存在动力学不稳定。
